% !PRINTED 247
\weberpage{247}
Man kann das System $\mathfrak{C}$ einen \gesperrt{endlichen Körper} nennen, da es nur eine endliche Anzahl von Zahlen enthält, die das charakteristische Merkmal eines Körpers, nämlich die unbeschränkte Ausführbarkeit der Rechenoperationen, ausgenommen die Division durch Null, aufweisen (Bd.\,I, \S.\,139)\textsuperscript{1)}. Wir wollen daher $\mathfrak{C}$ einen \gesperrt{Congruenzkörper} nennen.

Es soll $n$ der \gesperrt{Grad} und $p$ der \gesperrt{Modul} des Körpers $\mathfrak{C}$ heissen.

Ist $\alpha$ eine von Null verschiedene Zahl in $\mathfrak{C}$, so durchläuft das Product $\alpha\xi = \eta$ zugleich mit $\xi$ das volle Zahlensystem $\mathfrak{C}$. Schliessen wir $\xi = 0$ aus, so kommt auch $\eta = 0$ nicht vor, und das Product $\Pi$ aller $\xi$ stimmt mit dem Producte aller $\eta$ überein und ist von Null verschieden. Durch Multiplication aller Gleichungen $\alpha\xi = \eta$ folgt aber

\[
\alpha^{p^n-1}\Pi = \Pi,
\]

und nach Abwerfung des gemeinsamen Factors $\Pi$ ergiebt sich

\gesperrt{5. der Fermat'sche Satz:}

\weq{3}{\alpha^{p^n-1} = 1.}

Multiplicirt man mit $\alpha$, so erhält man diesen Satz in der Form

\weq{4}{\alpha^{p^n} = \alpha,}

in der er auch noch für $\alpha = 0$ besteht. Ist $\alpha$ von Null verschieden, so giebt es wegen (3) einen \gesperrt{kleinsten} positiven Exponenten $e$, für den

\weq{5}{\alpha^e = 1}

ist, und für den folglich die Potenzen $1, \alpha, \alpha^2, \ldots, \alpha^{e-1}$ von einander verschieden sind. Man sagt dann, $\alpha$ \gesperrt{gehört zum Exponenten} $e$. Dieser Exponent $e$ muss ein Theiler von $p^n - 1$ sein. Denn ist $\alpha^m = 1$ für irgend einen Exponenten $m$, so setzen wir $m = qe + e'$, worin $q$ eine ganze Zahl und $e' < e$ ist. Dann ist auch $\alpha^{e'} = 1$, und folglich muss $e' = 0$, d.\,h. $m$ durch $e$ theilbar sein. Unter diesen Exponenten $m$ findet sich auch $p^n - 1$. Setzen wir also

\[
p^n - 1 = ef,
\]

so wird $\alpha^h$ immer dann zum Exponenten $e$ gehören, wenn $h$ relativ prim zu $e$ ist. Denn dann und nur dann ist $he$ das kleinste positive, durch $e$ theilbare Vielfache von $h$.

\textsuperscript{1)} Man sehe des Verfassers Abhandlung: Die allgemeinen Grundlagen der \gesperrt{Galois'schen} Gleichungstheorie. Mathematische Annalen, Bd.\,43.
